\Askhsh{33596_SOLUTION}{1}
\begin{ylh}{1}
    \item [Ύλη από εικόνα]
\end{ylh}
\begin{wrapfigure}[24]{r}{8cm} % Αριθμός γραμμών και πλάτος ρυθμίζονται αναλόγως
    \centering
    \begin{tikzpicture}
    \begin{axis}[
        axis lines = middle,
        xmin = -0.3, xmax = 2.7,
        ymin = -0.5, ymax = 2.2,
        xtick = {-0.2, 0.2, 0.4, 0.6, 0.8, 1, 1.2, 1.4, 1.6, 1.8, 2, 2.2, 2.4, 2.6},
        ytick = {-0.4, -0.2, 0, 0.2, 0.4, 0.6, 0.8, 1, 1.2, 1.4, 1.6, 1.8, 2},
        grid = both,
        grid style = {gray!30},
        samples = 100,
        width = 7.5cm, height = 7.5cm, % Ρύθμιση μεγέθους
        tick label style = {font=\scriptsize},
        every axis x label/.style = {at={(current axis.right of origin)}, anchor=north west},
        every axis y label/.style = {at={(current axis.above origin)}, anchor=south},
        % Για σχεδίαση, χρησιμοποιούμε μια προσέγγιση για το x_o που ικανοποιεί x_o^2 + ln(x_o) - 2 = 0
        % x_o περίπου 1.3
        declare function={
            x_o_val = 1.30; % Τιμή του x_o, κοντά στη λύση του x^2 + ln(x) - 2 = 0
            ln_x_o_val = ln(x_o_val);
            x_B_val = x_o_val^3 - x_o_val; % x-συντεταγμένη του B
        },
    ]
        % Πράσινη καμπύλη: y = ln(x)
        \addplot[secondarycolor, thick, domain=0.05:2.6] {ln(x)}; % Το πεδίο ορισμού πρέπει να είναι > 0
        
        % Σημείο A (0,2)
        \node[above] at (axis cs:0,2) {A};
        \addplot[only marks, mark=*, mark size=1.5pt, black] coordinates {(0,2)};

        % Σημείο M (x_o, ln x_o)
        \node[below right] at (axis cs:x_o_val,ln_x_o_val) {M};
        \addplot[only marks, mark=*, mark size=1.5pt, black] coordinates {(x_o_val,ln_x_o_val)};

        % Σημείο B (x_B, 0)
        \node[below] at (axis cs:x_B_val,0) {B};
        \addplot[only marks, mark=*, mark size=1.5pt, black] coordinates {(x_B_val,0)};

        % Ευθεία AM (γκρι)
        \addplot[gray, thick] coordinates {(0,2) (x_o_val,ln_x_o_val)};

        % Εφαπτομένη ευθεία (ε) (κύριο χρώμα, ετικέτα)
        % Εξίσωση: y - ln(x_o) = (1/x_o)(x - x_o)
        \addplot[maincolor, thick, domain=x_B_val-0.1:2.6] {ln_x_o_val + (1/x_o_val)*(x - x_o_val)};
        \node[above left] at (axis cs:2.1, ln_x_o_val + (1/x_o_val)*(2.1 - x_o_val)) {$(\varepsilon)$};
        
        % Σύμβολο ορθής γωνίας στο M
        % Συντεταγμένες για M, A, και ένα σημείο στην εφαπτομένη (π.χ. B)
        \coordinate (M_pt) at (axis cs:x_o_val, ln_x_o_val);
        \coordinate (A_pt) at (axis cs:0, 2);
        \coordinate (B_pt_for_angle) at (axis cs:x_B_val, 0); % Χρησιμοποιούμε το B ως σημείο στην εφαπτομένη για τη σχεδίαση της γωνίας
        
        \draw[fill=secondarycolor!20, draw=black]
            ($(M_pt)!0.15cm!(A_pt)$) --
            ($(M_pt)!0.15cm!90:(A_pt)$) --
            ($(M_pt)!0.15cm!90:(B_pt_for_angle)$) --
            ($(M_pt)!0.15cm!(B_pt_for_angle)$) -- cycle;
    \end{axis}
    \end{tikzpicture}
\end{wrapfigure}
\begin{alist}
\item [α.] Η απόσταση ΑΚ συναρτήσει του $x > 0$ είναι
    \[ d(x) = \sqrt{(x-0)^2 + (\ln x - 2)^2} = \sqrt{x^2 + \ln^2 x - 4\ln x + 4} . \]
\item [β.] Η απόσταση $d(x)$ γίνεται ελάχιστη για $x=x_o$ και επειδή η $d(x)$ είναι παραγωγίσιμη συνάρτηση στο $(0,+\infty)$ από το θεώρημα του Fermat έχουμε ότι $d'(x_o)=0$. Όμως
    \[ d'(x) = \frac{2x + 2 \frac{\ln x}{x} - \frac{4}{x}}{2\sqrt{x^2 + \ln^2 x - 4\ln x + 4}} = \frac{x^2 + \ln x - 2}{x \cdot \sqrt{x^2 + \ln^2 x - 4\ln x + 4}} \]
    οπότε $d'(x_o)=0 \Rightarrow x_o^2 + \ln x_o - 2 = 0$ (1).
\item [γ.] Η εφαπτομένη της $C_f$ στο Μ έχει εξίσωση $(\varepsilon) : y - \ln x_o = \frac{1}{x_o}(x-x_o)$.
\begin{rlist}
\item [i.] Για να είναι κάθετη στην AM πρέπει και αρκεί
    $\lambda_{\varepsilon} \cdot \lambda_{AM} = -1 \Leftrightarrow \frac{1}{x_o} \cdot \frac{\ln x_o - 2}{x_o} = -1 \Leftrightarrow \ln x_o - 2 = -x_o^2 \Leftrightarrow x_o^2 + \ln x_o - 2 = 0 , που ισχύει από το β).$
\item [ii.] Για $y=0$ στην εξίσωση της $(\varepsilon)$ έχουμε
    $0 - \ln x_o = \frac{1}{x_o}(x-x_o) \Leftrightarrow -x_o \ln x_o = x-x_o \Leftrightarrow x = x_o - x_o \ln x_o (2).$
    Από την (1) έχουμε ότι $\ln x_o = 2 - x_o^2$ και έτσι η (2) γίνεται $x = x_o - x_o(2-x_o^2) \Leftrightarrow x = x_o^3 - x_o$. Συνεπώς η $(\varepsilon)$ τέμνει τον άξονα $xx'$ στο σημείο $B(x_o^3 - x_o, 0)$.
\end{rlist}
\end{alist}